% =====================================================================
%  The Onix Protocol — single-source LaTeX manuscript (Ledger template)
%  Author: Anatoly Piskunov (ORCID 0009-0000-8883-4111)
%
%  TARGET CLASS: ledger.cls  (LedgerJournal/Ledger-Latex-Template, GitHub)
%    Requires ledger.cls, ledgerbib.bst and an images/ folder on the build
%    path.  Compile with pdflatex.  This same .tex is the arXiv source (arXiv
%    primary cs.GT, cross-list q-fin.TR).
%
%  NOTES FOR THE FINAL BUILD
%    * ledger.cls already loads amsmath, amssymb, amsthm, hyperref, cleveref,
%      enumitem, tabularx, titling, titlesec, mathptmx, microtype.  The
%      preamble below therefore stays minimal to avoid clashes.
%    * References use a numbered thebibliography with parenthetical author-year
%      cites in prose (as in the source draft).  Ledger house style is endnotes
%      (\endnote/\theendnotes); converting is mechanical if the editor asks.
%    * No figures/images are used (the paper is math + tables only).
%
%  CHANGE OVER THE JULY 2026 DRAFT: this version weaves in the finalized F1
%  early-exit "deferred outcome-contingent claim" mechanism (new Section on
%  Early Exit + Theorem 2), reconciled with the leverage subsystem, governance
%  parameters, limitations and experimental hypotheses.  The mechanism is
%  live-verified on the VIZ testnet (full cycle: cancel -> claim recorded ->
%  settle paid at the capped bucket with a visible haircut -> claim consumed).
% =====================================================================
\documentclass{ledger}

% --- theorem environments (amsthm is loaded by the class) ---
\theoremstyle{plain}
\newtheorem{theorem}{Theorem}
\newtheorem{corollary}{Corollary}
\theoremstyle{remark}
\newtheorem*{remark}{Remark}

% --- Ledger front-matter counters / metadata (edit at typesetting) ---
\newcommand{\thefirstpagenum}{1}
\newcommand{\thelastpagenum}{1}
\newcommand{\theyear}{2026}
\newcommand{\thevol}{X}
\newcommand{\thedoi}{DOI 10.5195/LEDGER.\theyear.XXX}

\hypersetup{
  pdfauthor={Anatoly Piskunov},
  pdftitle={Combining Parimutuel Settlement with Automated Market Makers for Liquidity-Guaranteed Prediction Markets: The Onix Protocol}
}

% Clean the template banner: show only the title (remove "LEDGER LATEX TEMPLATE").
\renewcommand{\pretitle}{\begin{center}\fontsize{18pt}{22pt}\bfseries\selectfont}
\renewcommand{\posttitle}{\par\end{center}\vspace{0.6em}}

\title{Combining Parimutuel Settlement with Automated Market Makers for
Liquidity-Guaranteed Prediction Markets: The Onix Protocol}

\author{Anatoly Piskunov\thanks{Independent Researcher --- VIZ Distributed
Ledger. Email: \href{mailto:anatoly.piskunov@gmail.com}{anatoly.piskunov@gmail.com}.
ORCID: \href{https://orcid.org/0009-0000-8883-4111}{0009-0000-8883-4111}.}}

\date{2026}

\begin{document}
\maketitle

\begin{abstract}
Prediction markets aggregate dispersed information into probabilistic forecasts
that consistently outperform polls and expert panels, yet their adoption is
constrained by a structural liquidity problem: providers of market depth bear
adverse selection risk that deters retail participation. We present the
\textbf{Onix Protocol}, a hybrid architecture that decouples the \emph{pricing}
function from the \emph{settlement} function in prediction markets. By pairing
automated market maker pricing --- a Constant Product Market Maker (CPMM) for
binary outcomes and a Logarithmic Market Scoring Rule (LMSR) for multi-outcome
markets --- with parimutuel (totalizator) settlement, we achieve a
\textbf{structural guarantee} that liquidity-provider principal is never at risk
from betting outcomes. We formalize the protocol's economic invariants, prove
the LP principal guarantee for both market types, and resolve the central
tension of any curve-priced market --- that an early exit through the pricing
curve realizes a trading P\&L against the very LPs the protocol protects --- with
a \emph{deferred outcome-contingent claim} mechanism, for which we prove a
second no-mint / LP-safe invariant. We describe a ``Lazy'' liquidity pool
enabling passive retail participation, analyze the dispute resolution mechanism
under DAO governance, and discuss the experimental hypotheses this system is
designed to test. The protocol is implemented as consensus-level operations on
the VIZ distributed ledger.
\end{abstract}

\begin{keywords}
\item{prediction markets} \item{automated market makers} \item{parimutuel betting} \item{LMSR} \item{liquidity provision} \item{DAO governance} \item{mechanism design}
\end{keywords}

\clearpage
\tableofcontents
\clearpage

% =====================================================================
\section{Introduction}\label{sec:intro}

\subsection{The Prediction Market Liquidity Problem}

Prediction markets are among the most reliable mechanisms for aggregating
dispersed information into actionable probabilistic forecasts. Empirical
evidence from political forecasting (Berg, Forsythe, Nelson \& Rietz, 2008),
sports betting (Levitt, 2004), and internal corporate markets (Cowgill, Wolfers
\& Zitzewitz, 2009) demonstrates that market prices consistently outperform
alternative forecasting methods.

Despite this track record, prediction market deployment remains limited. The
fundamental constraint is not demand for forecasts but \emph{supply of
liquidity}. Every existing prediction market architecture requires liquidity
providers (LPs) to accept one of three risk profiles:

\begin{enumerate}
  \item \textbf{Impermanent loss} in AMM-based markets (Adams, Zinsmeister \&
  Robinson, 2020; Adams et al., 2021), where LPs systematically lose to informed
  traders;
  \item \textbf{Adverse selection} in limit-order-book markets, where market
  makers bear inventory risk against better-informed counterparties;
  \item \textbf{Bounded subsidy loss} in LMSR markets (Hanson, 2003), where the
  market maker risks up to $b \ln N$ against correctly-informed crowds.
\end{enumerate}

These risk profiles restrict liquidity provision to sophisticated actors with
active management capabilities, creating thin markets with high slippage that
deter bettor participation --- a self-reinforcing cycle:

\[
\text{thin markets} \rightarrow \text{high slippage} \rightarrow \text{poor UX}
\rightarrow \text{few bettors} \rightarrow \text{low fees} \rightarrow
\text{no LP incentive} \rightarrow \text{thin markets}
\]

\subsection{Research Questions}

This work poses and seeks to answer the following questions through a deployed
experimental system:

\textbf{Q1.} \emph{Can the pricing and settlement functions of a prediction
market be architecturally separated such that LP principal is structurally
insulated from betting outcomes?}

\textbf{Q2.} \emph{Does a parimutuel settlement mechanism, where losers
exclusively fund winners, eliminate the adverse selection problem inherent in
AMM-based and LMSR-based prediction markets?}

\textbf{Q3.} \emph{Can automated liquidity deployment via a pooled mechanism
(``Lazy Pool'') reduce the barrier to LP entry sufficiently to bootstrap market
depth without active market selection by individual LPs?}

\textbf{Q4.} \emph{Is a DAO-governed dispute resolution mechanism, with
stake-weighted public voting and bonded oracles, a viable alternative to
centralized arbitration or optimistic oracle schemes?}

\textbf{Q5.} \emph{What governance-token utility emerges from a prediction
market protocol where LP positions carry voting weight, and does this create a
sustainable token-demand feedback loop?}

\subsection{Contributions}

Our contributions are:

\begin{enumerate}
  \item \textbf{A formal proof} that combining AMM pricing with parimutuel
  settlement yields a structural LP principal guarantee for both binary (CPMM)
  and multi-outcome (LMSR) prediction markets (Theorem~\ref{thm:lp}).
  \item \textbf{A deferred outcome-contingent claim mechanism} that resolves the
  curve-exit / parimutuel tension --- allowing bettors and leveraged positions
  to exit early without realizing a trading P\&L against LP depth --- together
  with a proof (Theorem~\ref{thm:earlyexit}) that it cannot create an uncovered
  shortfall, mint tokens, or touch LP principal.
  \item \textbf{The ``Lazy Pool'' mechanism} --- an automated capital deployment
  system with graduated recall and opportunity-cost protection that enables
  passive retail LP participation, plus an opt-in, default-off pool-funded
  leverage subsystem whose pool exposure is \emph{bounded} and isolated from the
  Theorem~\ref{thm:lp} LP guarantee.
  \item \textbf{A two-mode dispute resolution framework} combining bonded
  oracles with DAO committee arbitration, including a full game-theoretic
  analysis of oracle incentive compatibility.
  \item \textbf{An experimental protocol} implemented as consensus-level
  operations on a production distributed ledger, with explicit hypotheses and
  measurable outcomes.
\end{enumerate}

% =====================================================================
\section{Related Work}\label{sec:related}

\subsection{Market Scoring Rules}

Hanson (2003) introduced the Logarithmic Market Scoring Rule (LMSR), in which a
market maker quotes prices via a softmax function over quantity parameters. The
LMSR guarantees a bounded loss to the market maker of $b \ln N$ and ensures
prices sum to unity. However, this bound is also the LP's maximum risk --- if
the crowd correctly predicts the outcome, the market maker loses the entire
subsidy. This has limited LMSR adoption to corporate environments (Microsoft,
Inkling) where the operator absorbs losses.

\subsection{Constant Product Market Makers}

The constant product formula $x \cdot y = k$, popularized by Uniswap (Adams,
Zinsmeister \& Robinson, 2020), provides continuous pricing for two-asset pairs.
When applied to binary prediction markets, the reserves represent opposing
outcomes and the invariant provides natural probability pricing. However,
standard AMM LPs suffer impermanent loss --- empirical analysis shows over 50\%
of Uniswap v3 LPs underperform buy-and-hold (Adams et al., 2021).

\subsection{Conditional Token Frameworks}

Polymarket employs the Gnosis Conditional Tokens Framework (CTF), where
positions are ERC-1155 tokens with split/merge operations to enforce price
coherence ($\sum p_i = 1$). This adds composability but introduces complexity
and requires an external arbitrage mechanism. UMA's Optimistic Oracle provides
dispute resolution through a challenge-response game with economic bonds.

\subsection{Parimutuel Betting}

The parimutuel (totalizator) model, originating from Pierre Oller's totalisator
(1865), pools all bets and distributes the losers' stakes proportionally among
winners. This guarantees the house never pays from its own capital. However,
traditional parimutuel systems lack real-time pricing --- odds are only final at
pool close --- making them unsuitable for continuous information aggregation.

\subsection{Our Approach}

The Onix Protocol combines the real-time pricing capability of AMMs/LMSR with
the zero-LP-risk property of parimutuel settlement. The AMM or LMSR serves
exclusively as a \emph{pricing engine} --- determining implied probabilities and
assigning weights --- while settlement follows the parimutuel model where losers
fund winners and LP principal is architecturally separated from the payout flow.

This positions Onix differently from two other hybrid directions in the
literature. First, \emph{liquidity-sensitive} market makers (Othman, Pennock,
Reeves \& Sandholm, 2013; Abernethy, Chen \& Vaughan, 2011) keep the operator as
counterparty but bound or shape its loss via the cost-function geometry; the
operator can still lose. Onix instead removes the counterparty role from the LP
altogether --- the curve only \emph{prices}, it never \emph{pays}. Second,
\emph{CLOB+AMM} hybrids (e.g. order books backstopped by an AMM, as explored in
several DEX designs and in Polymarket's CLOB with AMM-style fallbacks) blend two
pricing venues but inherit inventory/adverse-selection risk from the
market-making side. Onix keeps a single continuous pricing venue and moves the
risk-bearing entirely out of pricing and into the parimutuel pool. To our
knowledge, pairing a continuous AMM/LMSR pricer with parimutuel
(losers-fund-winners) settlement to obtain an unconditional LP-principal
guarantee has not previously been formalized.

% =====================================================================
\section{Protocol Model}\label{sec:model}

\subsection{Participants and Roles}

We define the following participant set $\mathcal{P}$:

\begin{center}
\begin{tabular}{lll}
\hline
\textbf{Role} & \textbf{Symbol} & \textbf{Function} \\
\hline
Market Creator & $c \in \mathcal{P}$ & Defines market parameters, provides seed liquidity \\
Oracle & $o \in \mathcal{P}$ & Registers insurance bond $I_o$, accepts markets, resolves outcomes \\
Bettor & $b_i \in \mathcal{P}$ & Places bets on outcomes, receives outcome tokens \\
Liquidity Provider & $\ell_j \in \mathcal{P}$ & Deposits capital into market pools or the Lazy Pool \\
Dispute Resolver & $\mathcal{D}$ & Committee (stake-weighted vote) or named account \\
\hline
\end{tabular}
\end{center}

\subsection{Market Definition}

A market $\mathcal{M}$ is a tuple:

\[
\mathcal{M} = (o, \tau, \mathbf{O}, t_{\text{bet}}, t_{\text{res}},
\boldsymbol{\theta}, L)
\]

where $o$ is the designated oracle, $\tau \in \{0, 1\}$ is the market type
(binary or multi), $\mathbf{O} = \{O_1, \ldots, O_N\}$ is the outcome set
($N = 2$ for binary, $3 \leq N \leq 10$ for multi), $t_{\text{bet}}$ and
$t_{\text{res}}$ are betting and resolution expiration times,
$\boldsymbol{\theta}$ is the fee parameter vector, and $L$ is the initial
liquidity.

\subsection{Fee Structure}\label{sec:fees}

All fees are expressed in basis points (bp), where $10000 \text{ bp} = 100\%$.
The fee vector is:

\[
\boldsymbol{\theta} = (\theta_{\text{oracle}}, \theta_{\text{creator}},
\theta_{\text{liq}})
\]

Fees are computed exclusively at resolution from the losing side's aggregate
stake $S_{\text{lose}}$ and are never deducted at bet placement time:

\begin{align*}
f_{\text{oracle}}  &= \lfloor S_{\text{lose}} \cdot \theta_{\text{oracle}} / 10000 \rfloor \\
f_{\text{creator}} &= \lfloor S_{\text{lose}} \cdot \theta_{\text{creator}} / 10000 \rfloor \\
f_{\text{liq}}     &= \lfloor S_{\text{lose}} \cdot \theta_{\text{liq}} / 10000 \rfloor \\
W &= S_{\text{lose}} - f_{\text{oracle}} - f_{\text{creator}} - f_{\text{liq}}
\end{align*}

where $W$ is the \textbf{winners' pool} available for distribution to winning
bettors.

\subsection{Market State Machine}\label{sec:fsm}

The lifecycle of $\mathcal{M}$ follows a finite state machine:

\[
q_0 \xrightarrow{\text{oracle accepts}} q_1 \xrightarrow{t \geq t_{\text{bet}}}
q_2 \xrightarrow{\text{oracle resolves}} q_3 \xrightarrow{\text{grace period}}
\text{paid}
\]

Two paths lead out of $q_0$ back to the terminal deleted state $q_{-1}$, both
returning the seed liquidity to the creator (the anti-spam creation fee, already
paid to the DAO fund at creation, is \emph{not} refunded): an explicit rejection
$q_0 \xrightarrow{\text{oracle rejects}} q_{-1}$, and a \emph{timeout}
$q_0 \xrightarrow{t \geq t_0 + \tau_{\text{accept}}} q_{-1}$ triggered by the
per-block cron when the oracle neither accepts nor rejects within the
governance-defined acceptance window $\tau_{\text{accept}}$ (default one hour).
The timeout bounds how long an unresponsive oracle can hold a creator's seed
capital hostage.

% =====================================================================
\section{Pricing Mechanisms}\label{sec:pricing}

\subsection{Binary Markets: Constant Product Market Maker}

For binary markets ($N = 2$), the protocol maintains two reserves $R_A$ and
$R_B$ subject to the invariant:

\[
k = R_A \cdot R_B
\]

\textbf{Initialization.} Given seed liquidity $L$:

\[
R_A = \lfloor L / 2 \rfloor, \quad R_B = L - R_A, \quad k = R_A \cdot R_B
\]

\textbf{Bet placement.} When bettor $b_i$ wagers amount $a$ on outcome $A$:

\begin{align*}
R'_A &= R_A + a \\
R'_B &= \lfloor k / R'_A \rfloor \\
w_i  &= R_B - R'_B \quad \text{(tokens received)}
\end{align*}

Here $w_i$ is the bettor's \emph{weight}. It is a \textbf{computed} quantity ---
the number of outcome tokens minted by the curve --- and is in general
\textbf{not} equal to the staked amount $a$: depending on the current reserves,
$w_i$ may be larger or smaller than $a$. The weight is used solely to apportion
the winners' pool at settlement (\Cref{sec:unified}); it is never a
currency-denominated claim.

\textbf{Implied probability.} The market-implied probability for each outcome is:

\[
P(A) = \frac{R_A}{R_A + R_B}, \quad P(B) = \frac{R_B}{R_A + R_B}
\]

Note that $P(A) + P(B) = 1$ by construction, eliminating the need for any
external arbitrage mechanism to enforce price coherence.

\subsection{Multi-Outcome Markets: LMSR}

For markets with $N > 2$ outcomes, the protocol uses the Logarithmic Market
Scoring Rule with softmax pricing.

\textbf{Cost function:}

\[
C(\mathbf{q}) = b \cdot \ln\left(\sum_{j=1}^{N} \exp(q_j / b)\right)
\]

where $\mathbf{q} = (q_1, \ldots, q_N)$ are quantity parameters and $b$ is the
liquidity parameter.

\textbf{Price function (softmax):}

\[
p_i(\mathbf{q}) = \frac{\exp(q_i / b)}{\sum_{j=1}^{N} \exp(q_j / b)}
\]

By the definition of softmax, $\sum_{i=1}^{N} p_i = 1$ identically.

\textbf{Cost of a trade.} The cost to purchase $\Delta$ tokens on outcome $i$:

\[
\text{cost}(\Delta, i) = C(\mathbf{q} + \Delta \cdot \mathbf{e}_i) - C(\mathbf{q})
\]

\textbf{Liquidity parameter.} The parameter $b$ is funded by the LP subsidy $S$:

\[
b = \frac{S}{\ln N}
\]

\textbf{Numerical stability.} The log-sum-exp computation uses the standard
identity:

\[
\ln\left(\sum_j \exp(x_j)\right) = \max(\mathbf{x}) +
\ln\left(\sum_j \exp(x_j - \max(\mathbf{x}))\right)
\]

% =====================================================================
\section{Parimutuel Settlement and LP Guarantee}\label{sec:settlement}

\subsection{Unified Settlement Model}\label{sec:unified}

Both market types share an identical settlement mechanism. At resolution, the
oracle declares winning outcome $O^* \in \mathbf{O}$. Let $\mathcal{W}$ be the
set of bets on $O^*$ and $\mathcal{L}$ be all other bets.

\[
S_{\text{lose}} = \sum_{b_i \in \mathcal{L}} a_i
\]

where $a_i$ is the bet amount. The winners' pool $W$ is computed as in
\Cref{sec:fees}.

Each winning bettor $b_i \in \mathcal{W}$ receives:

\[
\pi_i = a_i + \frac{w_i}{\sum_{b_j \in \mathcal{W}} w_j} \cdot W - \tau_i
\]

where $w_i$ is the bettor's weight (tokens from CPMM or LMSR) and $\tau_i$ is
the time penalty on profit (\Cref{sec:timepenalty}).

\subsection{Theorem: LP Principal Guarantee}\label{sec:thm1}

\begin{theorem}[LP Principal Guarantee]\label{thm:lp}
Under parimutuel settlement, the total money distributed equals the total money
received, with LP principal $L$ returned unconditionally. The LP principal is
never at risk from betting outcomes.
\end{theorem}

\begin{proof}
Consider the full money flow through a market:

\[
\text{Money}_{\text{IN}} = L + \sum_{b_i \in \mathcal{W}} a_i +
\sum_{b_i \in \mathcal{L}} a_i = L + \sum_{\text{all bets}} a_i
\]

where the first sum $\sum_{b_i \in \mathcal{W}} a_i$ is the aggregate principal
staked by \emph{winners} (returned to them in full at settlement) and
$S_{\text{lose}} = \sum_{b_i \in \mathcal{L}} a_i$ is the aggregate stake
forfeited by \emph{losers} (the sole source of the winners' pool and all fees).
At resolution, the outgoing flows are:

\[
\text{Money}_{\text{OUT}} = L + \sum_{b_i \in \mathcal{W}} a_i + W +
f_{\text{oracle}} + f_{\text{creator}} + f_{\text{liq}}
\]

Substituting $W = S_{\text{lose}} - f_{\text{oracle}} - f_{\text{creator}} -
f_{\text{liq}}$:

\begin{align*}
\text{Money}_{\text{OUT}}
&= L + \sum_{b_i \in \mathcal{W}} a_i + S_{\text{lose}} \\
&= L + \sum_{b_i \in \mathcal{W}} a_i + \sum_{b_i \in \mathcal{L}} a_i
 = L + \sum_{\text{all bets}} a_i = \text{Money}_{\text{IN}} \qquad \qedhere
\end{align*}
\end{proof}

\begin{corollary}
The LP principal guarantee holds regardless of the weight assignment mechanism,
the number of bettors, the distribution of bets across outcomes, or the fee
parameters. It is a property of the settlement architecture, not the pricing
curve.
\end{corollary}

\begin{corollary}
The maximum total payout to winners is $S_{\text{lose}}$ (the entire losing
pool), regardless of weights assigned by the AMM or LMSR. Therefore, winner
payouts can never draw from LP capital.
\end{corollary}

\begin{remark}[integer arithmetic and solvency under rounding]
Theorem~\ref{thm:lp} is stated over the rationals, but the protocol operates on
integers (milli-VIZ), and every division applies the floor operator
$\lfloor \cdot \rfloor$. We show the guarantee survives discretization. Each fee
and each winner share is rounded \emph{down}:

\[
f_\bullet = \lfloor S_{\text{lose}} \cdot \theta_\bullet / 10000 \rfloor, \qquad
\hat{\pi}_i^{\text{profit}} = \left\lfloor \frac{w_i}{\sum_{j} w_j} \cdot W
\right\rfloor
\]

Since $\lfloor x \rfloor \leq x$, the realized winners' distribution satisfies
$\sum_{i \in \mathcal{W}} \hat{\pi}_i^{\text{profit}} \leq W$, and likewise
$\sum_\bullet f_\bullet \leq S_{\text{lose}} \cdot (\theta_{\text{oracle}} +
\theta_{\text{creator}} + \theta_{\text{liq}})/10000 \leq S_{\text{lose}}$.
Therefore the realized outflow obeys

\[
\text{Money}_{\text{OUT}}^{\text{realized}} = L + \sum_{i \in \mathcal{W}} a_i +
\sum_{i \in \mathcal{W}} \hat{\pi}_i^{\text{profit}} + \sum_\bullet f_\bullet
\;\leq\; L + \sum_{\text{all bets}} a_i = \text{Money}_{\text{IN}}
\]

so the system is never insolvent. The non-negative residual

\[
\epsilon = S_{\text{lose}} - \sum_{i \in \mathcal{W}}
\hat{\pi}_i^{\text{profit}} - \sum_\bullet f_\bullet \;\geq\; 0
\]

is bounded by $\epsilon < |\mathcal{W}| + 3$ milli-VIZ (one sub-unit per floored
quantity) and is swept to the DAO fund as dust. Thus rounding can only
\emph{under}-distribute, never over-distribute: LP principal $L$ remains exactly
conserved and the inequality $\text{Money}_{\text{OUT}} \leq
\text{Money}_{\text{IN}}$ holds at the granularity of the ledger.
\end{remark}

\subsection{Time Penalty on Late Bets}\label{sec:timepenalty}

To mitigate information asymmetry from bets placed near expiration (when outcome
uncertainty is reduced), the protocol imposes a configurable time penalty
applied \emph{only to profit}, never to principal.

Let $\Delta t = t_{\text{bet}} - t_{\text{now}}$ be the time to expiration and
$\omega$ be the penalty window. If $\Delta t < \omega$:

\[
r = 1 - \frac{\Delta t}{\omega}
\]

\[
\tau_{\text{rate}} = \begin{cases} r & \text{linear} \\ r^2 &
\text{quadratic (default)} \end{cases}
\]

\[
\tau_i = \lfloor \pi_{\text{profit},i} \cdot \tau_{\text{rate}} \cdot
\tau_{\max} / 10^6 \rfloor
\]

where $\pi_{\text{profit},i}$ is the bettor's profit share (excluding principal)
and $\tau_{\max}$ is the maximum penalty in micro-units.

\textbf{Invariant:} $\pi_i \geq a_i$ for all winning bettors, since the penalty
is applied exclusively to the profit component.

\subsection{LP Time-Weighted Fee Distribution}\label{sec:lpfee}

LP fee shares are distributed proportionally to the product of deposit amount
and remaining time to expiration:

\[
\omega_j = \text{amount}_j \cdot \max(1, t_{\text{bet}} - t_{\text{deposit},j})
\]

\[
f_{\text{share},j} = \left\lfloor \frac{f_{\text{pool}} \cdot \omega_j}{\sum_m
\omega_m} \right\rfloor
\]

This creates a strong incentive for early liquidity provision. In a market of
duration $T$, an LP depositing at time $t = 0$ earns $T$ times more per unit
capital than one depositing at $t = T - 1$.

% =====================================================================
\section{``Lazy'' Liquidity Pool}\label{sec:lazypool}

\subsection{Motivation}

Individual LP provision requires active market evaluation and selection. For
retail participants, this creates an impractical knowledge and effort barrier.
The ``Lazy'' Pool (so named because it requires no active market selection from
depositors) solves this by accepting deposits and automatically allocating
capital to activated markets.

\subsection{Pool Mechanics}\label{sec:poolmech}

The Lazy Pool $\mathcal{L}_P$ is a singleton contract maintaining:
\begin{itemize}
  \item $B_{\text{free}}$: unallocated balance
  \item $B_{\text{alloc}}$: capital deployed to active markets
  \item $B_{\text{earned}}$: realized profits
  \item $S_{\text{total}}$: total outstanding shares
  \item $\rho$: cumulative reward-per-share accumulator
\end{itemize}

\textbf{Deposit.} When user $u$ deposits amount $a$:

\[
\text{shares}_u = \begin{cases} a & \text{if } S_{\text{total}} = 0 \\
a \cdot S_{\text{total}} / B_{\text{free}} & \text{otherwise} \end{cases}
\]

The deposit is locked for a governance-defined period $t_{\text{lock}}$.

\textbf{Auto-allocation.} On market activation ($q_0 \rightarrow q_1$):

\[
a_{\text{alloc}} = B_{\text{free}} \cdot \alpha / 100
\]

subject to $a_{\text{alloc}} \geq a_{\min}$ and $B_{\text{alloc}} +
a_{\text{alloc}} \leq B_{\text{total}} \cdot \alpha_{\max} / 100$, where $\alpha$
is the per-market allocation percentage and $\alpha_{\max}$ is the maximum total
allocation cap.

Allocation is further gated by a \textbf{reward floor}: the pool provides
capital to a market only if its LP fee $f_{\text{liq}}$ meets a governance
minimum $f_{\text{liq}}^{\min}$ (default $2\%$). A market that offers LPs less
than this earns nothing worth the pool's opportunity cost, so the pool allocates
zero and the market's only depth is the creator's own seed. This makes the LP
fee an explicit price the market pays for automated pool liquidity, and prevents
the pool from subsidizing markets configured to route all fee revenue away from
liquidity providers.

The allocation formula includes oracle quality adjustments:

\[
a_{\text{alloc}} = B_{\text{free}} \cdot \frac{\alpha}{100} \cdot
(1 - \beta)^{n_o} \cdot (1 - \gamma)^{f_o}
\]

where $n_o$ is the oracle's active market count, $f_o$ is the oracle's active
fault stamp count, and $\beta, \gamma$ are governance-defined decay factors.

\subsection{Reward Distribution (Lazy Accounting)}\label{sec:lazyacct}

Rewards are distributed using a single global accumulator $\rho$
(reward-per-share), following the MasterChef pattern (SushiSwap, 2020; Leshner \&
Hayes, 2019):

When a market resolves and the pool's LP position returns profit
$\pi_{\text{pool}}$:

\[
\rho \leftarrow \rho + \frac{\pi_{\text{pool}} \cdot \text{PRECISION}}{S_{\text{total}}}
\]

Any user $u$'s pending reward is computable in $O(1)$:

\[
\text{reward}_u = \text{pending}_u + \frac{\text{shares}_u \cdot (\rho -
\rho_{\text{snapshot},u})}{\text{PRECISION}}
\]

User records are updated only when that user acts (deposit or withdraw), giving
$O(1)$ amortized cost regardless of participant count.

\subsection{Opportunity-Cost Protection}\label{sec:oppcost}

The auto-allocation mechanism creates an attack vector: a malicious oracle could
create long-duration zero-volume markets to lock pool capital. Three graduated
defense mechanisms address this:

\textbf{Graduated Recall.} The market duration is divided into $n_{\text{steps}}$
intervals. At each checkpoint $k$, if cumulative volume $V_k$ satisfies:

\[
V_k < a_{\text{alloc}} \cdot \theta_{\text{recall}} / 100
\]

then a fraction $\delta_{\text{recall}}$ of the current allocation is recalled to
the pool. After $n$ consecutive low-volume checkpoints, the retained allocation
is:

\[
a_{\text{retained}} = a_{\text{alloc}} \cdot (1 - \delta_{\text{recall}})^n
\]

For the default parameters ($\delta_{\text{recall}} = 10\%$, $n_{\text{steps}} =
10$), a completely idle market retains $(0.9)^{10} \approx 35\%$ of its original
allocation.

\textbf{Active Market Penalty.} Each active market from oracle $o$
multiplicatively reduces new allocations by factor $(1 - \beta)$. With $\beta =
5\%$ and 10 active markets, new allocations are $(0.95)^{10} \approx 60\%$ of the
base rate.

\textbf{Fault Stamps.} Bad outcomes (no-contest, missed deadlines, dispute
losses, zero-volume resolutions) generate penalty stamps on the oracle that
further reduce allocations by factor $(1 - \gamma)^{f_o}$. Stamps auto-expire
after a governance-defined clean-operation window.

\subsection{Emergency Withdrawal}

Users may exit the pool before lock expiration, subject to a penalty on
\emph{locked profit only}:

\[
\text{penalty} = \max(0, V_{\text{total}} - P_{\text{deposited}}) \cdot
\frac{S_{\text{locked}}}{S_{\text{total}}} \cdot \theta_{\text{emergency}} / 100
\]

where $V_{\text{total}}$ is the user's total value (shares + pending rewards),
$P_{\text{deposited}}$ is cumulative principal deposited, and $S_{\text{locked}}
/ S_{\text{total}}$ is the fraction of shares still locked. The penalty is
redistributed to remaining pool participants via $\rho$.

\subsection{Opt-In Leverage Funded by the Pool}\label{sec:leverage}

The Lazy Pool plays a second, optional role from the same $B_{\text{free}}$: it
funds a leverage subsystem. This subsystem is \textbf{opt-in, default-off, and
governed by a median kill-switch}. We treat it separately from the rest of the
paper for an important reason. \textbf{It is the one place where pool capital
takes on credit risk, so it is \emph{not} covered by the unconditional guarantee
of Theorem~\ref{thm:lp}.} Theorem~\ref{thm:lp} concerns the pool acting as a
\emph{market liquidity provider}, where principal is structurally insulated from
betting outcomes. Leverage instead has the pool act as a \emph{lender}, and
lending carries risk that we bound rather than eliminate.

\textbf{Position opening.} A bettor posts collateral $m$. The pool lends margin
$\lambda m$ at leverage factor $\lambda$ from $B_{\text{free}}$, subject to caps
on per-position size, pool fund fraction, and minimum market liquidity. No token
is minted: the total stake $(1+\lambda)m$ enters the market exactly as an
ordinary bet would. The position's curve weight, its collateral, and the lent
amount are recorded on a dedicated position object. The pool's obligation on the
position is

\[
\Omega = \lambda m \,(1 + r),
\]

where $r$ is the accrued interest. The loan plus interest is what the pool seeks
to recover.

\textbf{Liquidation against the recorded reserve snapshot.} The hard case in
leveraged binary markets is \emph{jump risk}: a discontinuous price move can
leave a naive ``sell at the current price'' liquidation unable to cover the
loan. Onix does not liquidate at the current price. It closes the position via
reverse-CPMM against the \textbf{reserve state recorded for that position}, so
the position's own market impact is unwound first rather than being paid for
twice. Two paths drive recovery:

\begin{enumerate}
  \item \textbf{Opposing-bet cascade} (\texttt{pm\_place\_bet}, liquidation
  reason 0) --- an incoming opposing bet triggers a forced close of the
  leveraged position against its recorded snapshot.
  \item \textbf{Settlement force-close} (\texttt{pm\_leverage\_resolve}) --- at
  resolution the position is closed and the loan repaid from its weight before
  any profit accrues to the borrower.
\end{enumerate}

On these paths the design target is full loan recovery, $\text{recovered} =
\min(V_{\text{close}}, \Omega) \ge \lambda m$, with the realized loan-and-interest
returning to the pool. We state this as a \textbf{design property, not a
theorem}: the leverage settlement math is implemented and exercised in
\texttt{consensus\_sim}, but a closed-form proof of recovery under all reserve
trajectories is left to future work (\Cref{sec:limitations}).

\textbf{Where the borrower's upside now goes.} Because leverage closes against
the curve rather than as a parimutuel share, its exit is governed by the
early-exit deferred-claim mechanism (\Cref{sec:earlyexit}): the pool takes its
obligation $\Omega$ from the close value first, and any residual $\max(V_{\text{close}}
- \Omega, 0)$ becomes an \emph{outcome-contingent deferred claim} rather than an
immediate payout. Leverage is thus a \textbf{leveraged directional bet}, not a
volatility harvest --- the borrower profits only if the chosen outcome wins and
the capped bucket has room. This is precisely what closes the historical
uncovered-shortfall path (\Cref{sec:earlyexit}); with it, the borrower's upside
can never draw on LP principal.

\textbf{Kill-switch decoupled from protection.} The kill-switch disables
\emph{new} openings only. The liquidation paths are deliberately \textbf{not}
gated by it, so disabling leverage never strips protection from already-open
positions.

\textbf{Risk isolation and reward.} Leverage lending is confined to a governed
fraction of $B_{\text{free}}$ (\texttt{leverage\_fund\_percent}), so worst-case
pool exposure is capped at the subsystem level and cannot reach the market-LP
principal that Theorem~\ref{thm:lp} protects. On the same-side voluntary-cancel
path a borrower closing into an adverse move can still leave the pool short of
$\Omega$; this shortfall is \textbf{bounded by the borrower's collateral $m$},
which the protocol holds as first-loss margin, and cannot exceed it. Interest
earned accrues to the pool through the same $\rho$ accumulator
(\Cref{sec:lazyacct}). Pool depositors therefore earn from two sources ---
losers'-pool fee shares (risk-free, Theorem~\ref{thm:lp}) and leverage interest
(bounded credit risk, opt-in) --- and the two are accounted identically but
governed independently.

\subsection{Design Decision: Real Depth Only (No Virtual/Phantom Liquidity)}\label{sec:realdepth}

A natural proposal is to seed a market's pricing curve with \emph{virtual}
(phantom) liquidity --- a reserve offset $\phi$ added to flatten price impact but
backed by no real capital and deleted at settlement, optionally tuned by
governance. In the closed bet$\rightarrow$cancel$\rightarrow$settle loop this is
value-conservative (it is the virtual-AMM technique), and it is tempting as a
cold-start ``stabilizer'' for new, thin markets. Onix \textbf{deliberately does
not implement it.} The same cold-start benefit is already delivered by Lazy-Pool
auto-allocation (\Cref{sec:poolmech}) --- but with \emph{real} capital, which
additionally earns fees, has an accountable owner, and follows demand per
market. We reject phantom depth because, applied carelessly, it damages the two
things the protocol exists to protect --- market structure and trust:

\begin{enumerate}
  \item \textbf{Forgeable depth erodes trust.} Depth is meaningful as a signal
  only because it is costly real capital at risk. Free virtual depth makes
  ``this market is deep and liquid'' a forgeable claim --- a thin or manipulated
  market can be dressed to look deep. Real capital makes the claim unforgeable.
  \item \textbf{It silently distorts information aggregation.} Virtual depth
  flattens the weight curve, weakening the reward for early, correct information
  and making the displayed price unresponsive to news (``stable because
  unmovable'' $=$ a stale forecast). The right magnitude is per-market and
  volume-dependent; a single governance constant cannot track it and, set too
  high, degrades the very price-discovery property a prediction market provides.
  \item \textbf{It is solvent only while never redeemed or used as collateral.}
  The moment depth backs a real outflow --- cancellations, early withdrawals,
  leverage loans, shared/cross-market pools --- the virtual part must be excluded
  everywhere or it leaks real money (e.g. a leverage loan sized or recovered
  against fake depth becomes real bad debt to pool depositors). Every ``size
  against liquidity / pay the LP'' branch becomes a footgun requiring ``\ldots but
  not the phantom part.'' Real capital removes this entire class of
  excludability bugs by construction.
  \item \textbf{It has no owner, no yield, no accountability.} Virtual depth
  bears no risk and earns no fee for anyone real; it is a service nobody is paid
  for and nobody answers for. For the markets it touches, it deletes the retail
  safe-yield product that is the protocol's core hypothesis (Q3, H1).
\end{enumerate}

The Lazy Pool is the same idea done with real numbers: auto-allocation smooths
the launch of new markets, but the capital is redeemable, leverage-safe,
fee-earning, owned by depositors, and self-correcting per market via graduated
recall (\Cref{sec:oppcost}). Onix therefore keeps \textbf{only real numbers} ---
every unit of depth is real capital that can be withdrawn, earns, and is
accountable. This is a conscious tradeoff: we forgo a cheap virtual stabilizer
to preserve the integrity of the price signal and the solvency of every
real-money path.

% =====================================================================
\section{Early Exit via Deferred Outcome-Contingent Claims}\label{sec:earlyexit}

\subsection{The Curve-Exit / Parimutuel Tension}\label{sec:exittension}

The Onix market is a hybrid: a CPMM (binary) / LMSR (multi) \textbf{curve} prices
both entry and early exit, while positions \emph{held to resolution} settle
parimutuel (\Cref{sec:settlement}). This creates a structural tension. Any
round-trip through the curve --- buy, then sell before settlement --- realizes a
trading P\&L against the curve's depth, which is the LPs, exactly like Uniswap
impermanent loss. But Theorem~\ref{thm:lp} promises LPs \emph{principal
protection} with fee-only upside. The two properties are in direct conflict on
any early-exit path.

Two consensus paths exit against the curve on a binary market:

\begin{itemize}
  \item \textbf{Leverage} (\texttt{liquidate\_position} / \texttt{pm\_leverage\_close})
  --- always, and force-closed at settlement, amplified by the pool loan
  (\Cref{sec:leverage}).
  \item \textbf{Regular bet cancel} (\texttt{cancel\_bet}, a curve-priced refund)
  --- during the betting window.
\end{itemize}

In the naive design both paths route a signed residual
$\text{residual} = \text{stake} - \text{curve\_refund}$ to the market's
$\text{forfeit\_pool}$. When an early exit is \emph{profitable}
($\text{curve\_refund} > \text{stake}$), the residual is negative and
$\text{forfeit\_pool}$ turns negative. At settlement the winners' pool is

\[
\text{winners\_pool} = S_{\text{lose}} - \text{fees} + \text{forfeit\_pool},
\]

so if early-exit (especially leveraged) profit outran the losing stakes,
$\text{winners\_pool} < 0$. Floored to zero, the shortfall --- call it
$\text{uncovered}$ --- is charged to LP principal, or, once LP principal is
exhausted, effectively \emph{minted}. This is not hypothetical: it is reachable
under a gate-respecting CPMM simulation (a one-sided pump produces
$\text{uncovered} = 7316$ milli-VIZ), and an adversarial replay corpus hit it in
$1163/1988$ paired trajectories.

The root cause is precise: \textbf{a curve-priced exit pays a bonding-curve value
that is not bounded by the losing pool}, whereas settlement pays parimutuel
(bounded by $S_{\text{lose}}$). The unbounded gap lands on the LP --- breaking
the very guarantee the architecture exists to provide.

\subsection{The Deferred-Claim Mechanism}\label{sec:defclaim}

Onix resolves the tension by making early exits \emph{not} extract curve value
from LPs at all. Instead, an exit records an \textbf{outcome-contingent deferred
claim}, funded at settlement from a \emph{bounded slice of the losing pool}.

\textbf{Recorded on exit.} A deferred-claim object stores
$\big(\text{position\_id},\ \text{kind} \in \{\text{bet}, \text{leverage}\},\
\text{chosen\_outcome},\ \text{claim\_amount},\ \text{exit\_time}\big)$.

\textbf{Regular bet cancel.} Principal is returned \emph{immediately and
unconditionally},
\[
\text{refund} = \min(\text{curve\_refund}, \text{stake}),
\]
because it is the bettor's own money. A cancel can cut losses or break even but
never realizes curve-profit at cancel time. The profit tail
\[
\text{claim\_amount} = \max(\text{curve\_refund} - \text{stake}, 0)
\]
becomes a deferred claim on the chosen outcome. Crucially, only a
\emph{non-negative} residual $\text{stake} - \text{refund} \ge 0$ is ever routed
to $\text{forfeit\_pool}$; the negative-residual path that produced
$\text{uncovered}$ is eliminated.

\textbf{Leverage close / liquidate.} Collateral is \emph{not} separately
returned --- it is first-loss margin for the pool. The pool recovers its
obligation $\Omega = \lambda m (1 + r)$ from the close value $V_{\text{close}}$
first; if $V_{\text{close}} < \Omega$, the collateral covers the gap. The
residual
\[
\text{claim\_amount} = \max(V_{\text{close}} - \Omega, 0)
\]
becomes a deferred claim on the chosen outcome. Two distinct predicates now
govern a leveraged position, and they must not be conflated: \emph{solvency}
($V_{\text{close}} \ge \Omega$, which governs loan recovery to the pool) versus
\emph{outcome-win} (which governs the right to a claim). A position can be solvent
yet on the losing outcome --- the pool is made whole and the claim is zero.

\subsection{Settlement Distribution}\label{sec:defsettle}

At settlement the deferred claims are funded from a capped bucket:

\begin{enumerate}
  \item Compute the bucket, clamped to the settlement headroom $H$
  (\Cref{sec:defthm}) so it can never exceed the pool available to winners:
  \[
  \text{bucket} = \min\!\left(\left\lfloor
  \frac{\texttt{pm\_early\_exit\_reward\_cap\_percent} \cdot S_{\text{lose}}}{10000}
  \right\rfloor,\; H\right)
  \]
  (default cap $=3300$ bp $=33\%$ of the losing pool; the clamp binds only when
  fees are large enough that the cap would otherwise over-draw the pool).
  \item Collect deferred claims on the \textbf{winning outcome only};
  losing-outcome claims pay zero.
  \item Pay them \textbf{FIFO by \text{exit\_time}} (first out, first paid) until
  the bucket is drained. There is no per-position cap --- FIFO ordering plus the
  total bucket \emph{is} the bound. A claim the remaining bucket cannot fully
  fund is paid partially; the unfunded remainder is a haircut.
  \item \textbf{Any unused bucket returns to the winners' pool}, where held
  winning bets share it parimutuel.
\end{enumerate}

\subsection{Theorem: Bounded Early Exit, No Uncovered Shortfall}\label{sec:defthm}

\begin{theorem}[Bounded Early Exit]\label{thm:earlyexit}
Let $c = \texttt{pm\_early\_exit\_reward\_cap\_percent}/10000 \in [0,1)$, and let
the funding bucket be clamped to the settlement headroom:
\[
\text{bucket} = \min\!\big(\lfloor c\, S_{\text{lose}}\rfloor,\; H\big), \qquad
H = \max\!\big(0,\; S_{\text{lose}} - \text{fees} - f_{\text{fixed}}^{\text{paid}}
+ \text{honest\_forfeits}\big),
\]
where $H$ is the winners' pool available immediately \emph{before} claims are paid
and $\text{honest\_forfeits} \ge 0$. Then, for \emph{any} market and \emph{any}
valid fee parameters, the paid early-exit claims never exceed the bucket, the
winners' pool is non-negative, and LP principal is never touched by early exits:
\[
\text{paid\_claims} \;\le\; \text{bucket} \;\le\; H, \qquad
\text{winners\_pool} \;=\; H - \text{paid\_claims} \;\ge\; 0.
\]
\end{theorem}

\begin{proof}
The FIFO distribution (\Cref{sec:defsettle}) stops paying once cumulative
payments reach the bucket, so $\text{paid\_claims} \le \text{bucket}$. By the
clamp, $\text{bucket} \le H$, hence $\text{paid\_claims} \le H$. The winners'
pool is exactly $H$ minus what was paid to early exiters, so
$\text{winners\_pool} = H - \text{paid\_claims} \ge 0$. No shortfall is charged
to LP principal, and because outflow never exceeds $S_{\text{lose}} + L +
\sum_{\mathcal{W}} a_i$, no tokens are minted. After the mechanism of
\Cref{sec:defclaim} only non-negative residuals reach $\text{forfeit\_pool}$, so
$\text{honest\_forfeits} \ge 0$ and $H$ is well defined.
\end{proof}

The clamp is what makes the guarantee \emph{unconditional}. Without it, one would
need the governance invariant $\theta_{\text{total}}/10000 + c \le 1$ (with
$\theta_{\text{total}} = \theta_{\text{oracle}} + \theta_{\text{creator}} +
\theta_{\text{liq}}$), giving the intuitive bound $\text{winners\_pool} \ge
(1-c)\,S_{\text{lose}} - \text{fees} \ge 0$. That invariant is \emph{not} enforced
by parameter validation --- per-market fees are bounded only by
$\theta_{\text{total}} \le 100\%$ at market creation, with no chain-level cap on
the creator or liquidity fee --- so a valid market with high fees combined with a
non-zero cap could otherwise drive $\text{winners\_pool}$ negative and charge the
shortfall to LP principal. In the default regime (fees $\approx 2\%$, cap $33\%$)
the $\min$ is inactive and $\text{bucket} = \lfloor c\,S_{\text{lose}}\rfloor$,
recovering the intended cap exactly; the clamp binds only in the adversarial
high-fee corner, where it preserves solvency at the cost of a deeper early-exit
haircut.

\begin{corollary}[No-mint, LP-safe, loser-never-profits]
Under Theorem~\ref{thm:earlyexit}: (i) an $\text{uncovered}$ shortfall is
impossible by construction, so no minting occurs and LP principal --- including
the Lazy Pool's market-LP principal --- is never touched by early exits or by
leverage impermanent loss; (ii) a claim on a \emph{losing} outcome pays zero, so
a losing outcome never profits; and (iii) held winners receive at least
$(1-c)\,S_{\text{lose}}$ of the losing pool plus any unused bucket.
\end{corollary}

\subsection{Economic Interpretation: A Bounded Liquidity Discount}\label{sec:defecon}

The deferred claim reframes early exit as a \emph{choice} with clean incentives.
A holder who stays to resolution receives a full parimutuel share of the losing
pool. A holder who exits early instead receives (a) immediate return of
principal (regular bet) or loan-satisfied collateral treatment (leverage), plus
(b) a \emph{capped, priority-ordered, outcome-contingent} claim of at most $c$ of
the losing pool, paid FIFO. Because the early-exit payoff is strictly a
capped-and-contingent version of the hold payoff, \textbf{there is no arbitrage
against holding}: one cannot exit early to extract \emph{more} than holding would
have yielded. The discount is deliberate --- it is the price of liquidity and of
transferring outcome risk earlier --- and it is exactly what makes the LP
guarantee survive a curve that also prices exits.

This also gives leverage a coherent economic identity. Under the deferred-claim
rule, borrowing from the Lazy Pool to open a position is a leveraged
\emph{directional} bet: the borrower amplifies exposure to a chosen outcome, the
pool is repaid its loan and interest first from the curve close value, and the
borrower's upside is a bounded contingent claim that competes FIFO with other
early exiters for the capped bucket. The volatility-harvest interpretation ---
open, ride the curve, cash out the swing at LP expense --- is structurally
removed.

\subsection{Governance Parameter and Implementation}\label{sec:defimpl}

The cap is a median-voted validator parameter
\texttt{pm\_early\_exit\_reward\_cap\_percent} (uint16, basis points, default
$3300$), added to \texttt{chain\_properties\_pm} with $\texttt{validate()}$
bounding it to $\le 10000$. Deferred claims live in a dedicated
\texttt{pm\_deferred\_claim\_object} (indexed by market and by exit order),
included in the snapshot allowlist with an import handler. A virtual operation
\texttt{pm\_early\_exit\_claim\_paid}$(\text{account}, \text{market},
\text{kind}, \text{outcome}, \text{claimed}, \text{paid})$ is emitted in the
settlement distribution loop --- exposing any bucket-exhaustion haircut via the
gap between \texttt{claimed} and \texttt{paid} --- and routed to the exiter's
account history (a bare balance adjustment would leave no trace). A read API
\texttt{get\_deferred\_claims} returns the FIFO queue, empty on a settled market
once claims are consumed.

Three defensive measures harden the mechanism against adversarial parameter
extremes and denial-of-service. First and most important, the funding bucket is
\textbf{clamped to the settlement headroom} $H$ (\Cref{sec:defthm}): the amount
actually drawn for claims is $\min(\lfloor c\,S_{\text{lose}}\rfloor, H)$, which is
what makes Theorem~\ref{thm:earlyexit} unconditional rather than contingent on a
$\text{fees} + \text{cap} \le 100\%$ governance invariant --- the latter is not
enforced by parameter validation, since the creator and liquidity fees have no
chain-level cap. Second, the per-market count of deferred-claim objects is capped
by \texttt{MAX\_PM\_DEFERRED\_CLAIMS\_PER\_MARKET}; beyond the cap an early exit
skips claim creation and the tail simply stays in the curve paying zero --- it is
deliberately \emph{not} routed to $\text{forfeit\_pool}$, which would double-count
and thereby mint. Third, an always-on diagnostic fires if $\text{uncovered} > 0$
ever occurs (unreachable once the bucket is clamped), and a $\lfloor \cdot
\rfloor_{\ge 0}$ floor on the winners' pool remains as a final backstop. These
are belt-and-suspenders: the clamp gives the guarantee by construction, and the
diagnostic makes any latent violation loud rather than silent.

\textbf{Verification.} The conservation identity of Theorem~\ref{thm:earlyexit}
was validated in a settlement model across adversarial scenarios (one-sided
pumps, thin losing sides, bucket exhaustion, fee extremes, and
$\Omega > \text{total stake}$ from long-lived leverage) and in the
\texttt{consensus\_sim} harness. It was further confirmed end-to-end on the VIZ
testnet: a cancel recorded a deferred claim; at settlement the virtual operation
reported the claim paid at the capped bucket with a visible haircut
($\texttt{claimed} > \texttt{paid}$), and the claim was consumed --- the full
lifecycle behaving exactly as the model predicts.

% =====================================================================
\section{Dispute Resolution Under DAO Governance}\label{sec:dispute}

\subsection{Bonded Oracle Model}

Each oracle $o$ must maintain an insurance bond $I_o \geq I_{\min}$. The bond
creates accountability: oracles who misresolve markets, miss deadlines, or lose
disputes have their insurance slashed. Oracle revenue derives from:

\begin{enumerate}
  \item A fixed per-market fee $f_{\text{fixed}}$ (compensating insurance
  staking)
  \item A percentage fee $f_{\text{oracle}}$ from the losers' pool at resolution
\end{enumerate}

These two fees enter the system at different points and must not be conflated.
The \textbf{percentage fee} $f_{\text{oracle}}$ is deducted from
$S_{\text{lose}}$ at resolution and appears in the money flow of
Theorem~\ref{thm:lp} (\Cref{sec:thm1}). The \textbf{fixed fee}
$f_{\text{fixed}}$ is a direct creator$\rightarrow$oracle transfer executed at
market \emph{acceptance}, before any betting; it is therefore \emph{not} part of
the betting/LP money flow --- it neither adds to the seed liquidity $L$ nor draws
from $S_{\text{lose}}$, and so does not appear in Theorem~\ref{thm:lp}. (It is
omitted from the fee vector $\boldsymbol{\theta}$ of \Cref{sec:fees} for the same
reason: $\boldsymbol{\theta}$ collects only the resolution-time, losers-funded
percentage fees.) For self-oracle markets no transfer occurs and
$f_{\text{fixed}} = 0$.

The oracle acceptance flow implements an offer-quote mechanism: the creator
publishes fee ceilings, and the oracle freezes its actual terms (bounded by both
the creator ceiling and a governance cap) at acceptance. Acceptance is
time-bounded by the window $\tau_{\text{accept}}$ (\Cref{sec:fsm}): a market left
pending beyond it is auto-expired by the per-block cron, which refunds the
creator's seed liquidity while retaining the anti-spam creation fee --- so an
unresponsive oracle cannot indefinitely freeze the seed.

\subsection{Two-Mode Dispute System}\label{sec:disputemodes}

Any bettor may challenge a resolution within a grace period
$\Delta t_{\text{grace}}$ by escrowing a dispute fee $d$. During disputes, all
payouts are frozen. The protocol supports two dispute modes per market:

\textbf{Committee mode ($\delta_{\text{mode}} = 0$).} The entire token-holder
electorate resolves the dispute via stake-weighted public vote. Each voter $v$'s
weight is:

\[
w_v = s_v + \text{shares}_{v,\text{pool}} \cdot \text{NAV}_{\text{pool}} /
S_{\text{total}}
\]

where $s_v$ is the voter's vesting shares and the second term converts Lazy Pool
stake into equivalent governance weight. Votes are public and revisable until
the voting period closes --- this is a deliberate design choice: disputes are
\emph{transparent public hearings}, not secret ballots. New evidence can change
votes.

This transparency carries a known cost. Because the running tally is visible, a
voter can defer until late in the window and condition its ballot on others'
revealed positions (a last-mover advantage). The protocol accepts this tradeoff
on purpose: for a DAO, the auditability and legitimacy of an open hearing are
judged more valuable than ballot secrecy, and a flipped outcome still requires
moving a stake-weighted majority. Manipulation is deterred not by secrecy but by
(i) the stake-weighted approval threshold $\theta_{\text{approve}}$, (ii) a
multi-day voting window that dilutes any single timing edge, and (iii) the
resolver-reward alignment analyzed in \Cref{sec:gametheory}. A commit-reveal
ballot was explicitly considered and \textbf{rejected} for this reason ---
concealing votes would defeat the public-hearing property that gives the verdict
its legitimacy. We revisit this as an honest limitation in \Cref{sec:limitations}.

The dispute is upheld if the approval percentage exceeds a governance threshold
$\theta_{\text{approve}}$:

\[
\frac{\sum_{v: \text{vote}_v = \text{approve}} w_v}{\sum_{v} w_v} \geq
\frac{\theta_{\text{approve}}}{10000}
\]

\textbf{Account mode ($\delta_{\text{mode}} = 1$).} A named dispute resolver
(recommended: multisig) issues a binding verdict.

\subsection{Dispute Outcomes and Incentive Compatibility}

\textbf{Oracle wrong (dispute upheld):}

\[
\text{reward}_{\text{pool}} = \min(d \cdot \mu, I_o)
\]

where $\mu$ is the reward multiplier. The reward pool is split:

\begin{itemize}
  \item Disputer receives $d$ (fee refund) $+\ \text{reward}_{\text{pool}} / \mu$
  \item Resolver/voters receive $\text{reward}_{\text{pool}} -
  \text{reward}_{\text{pool}} / \mu$
  \item Remaining insurance: additional penalty $\rightarrow$ DAO fund
\end{itemize}

\textbf{Oracle right (dispute rejected):}

\[
d \rightarrow \text{50\% resolver, 50\% oracle}
\]

\textbf{Oracle non-responsive (auto-close at 14 days):}

\begin{itemize}
  \item Disputer: fee refunded
  \item Oracle: $d$ slashed from insurance
  \item All bets and LP positions: full refund
  \item Slashed amount distributed proportionally to all participants
\end{itemize}

\subsection{Game-Theoretic Analysis}\label{sec:gametheory}

\textbf{Oracle incentive compatibility.} The oracle's expected payoff from
honest resolution vs. manipulation is:

\[
\mathbb{E}[\pi_{\text{honest}}] = f_{\text{fixed}} + f_{\text{oracle}} \quad
\text{(per market)}
\]

\[
\mathbb{E}[\pi_{\text{dishonest}}] = p_{\text{detect}} \cdot (-I_o \cdot
\theta_{\text{penalty}} - B_{\text{ban}}) + (1 - p_{\text{detect}}) \cdot
g_{\text{manipulation}}
\]

where $p_{\text{detect}}$ is the probability of dispute, $\theta_{\text{penalty}}$
is the insurance slash fraction, $B_{\text{ban}}$ is the present value of a ban
(lost future revenue), and $g_{\text{manipulation}}$ is the one-time
manipulation gain.

Honest behavior is a Nash equilibrium when:

\[
f_{\text{fixed}} + f_{\text{oracle}} > (1 - p_{\text{detect}}) \cdot
g_{\text{manipulation}} - p_{\text{detect}} \cdot (I_o \cdot \theta_{\text{penalty}}
+ B_{\text{ban}})
\]

The protocol ensures this by requiring $I_o \gg g_{\text{manipulation}}$ and
maintaining high $p_{\text{detect}}$ through public dispute visibility and
committee participation incentives.

\textbf{Disputer incentive.} A rational bettor files a dispute when the expected
reward exceeds the dispute fee:

\[
\mathbb{E}[\text{reward}_{\text{dispute}}] = p_{\text{upheld}} \cdot (d + d \cdot
(\mu - 1) / \mu) > d
\]

This simplifies to $p_{\text{upheld}} > 1/\mu$, establishing a natural threshold
for dispute filing.

\textbf{Committee voting incentive.} In committee mode, voters participate
because their Lazy Pool stake (and thus governance weight) directly earns a share
of the resolver reward pool. The expected reward for voter $v$ is:

\[
\mathbb{E}[\pi_v] = \frac{w_v}{\sum w_j} \cdot \text{voter\_reward\_pool} \cdot
p_{\text{upheld}}
\]

\subsection{No-Contest Declaration}

An oracle unable to verify an outcome may declare no-contest at reduced cost
($50\%$ of the dispute penalty from insurance). This creates a three-tier
incentive gradient:

\begin{center}
\begin{tabular}{lll}
\hline
\textbf{Action} & \textbf{Oracle Cost} & \textbf{Ban Risk} \\
\hline
Voluntary no-contest & $0.5 \cdot d$ from insurance & None \\
Dispute loss & $I_o \cdot \theta_{\text{penalty}} + \text{extra}$ & Yes \\
Missed deadline & $I_o \cdot \theta_{\text{miss}}$ & None \\
\hline
\end{tabular}
\end{center}

A no-contest declaration is itself disputable, with the resolver choosing from
three outcomes (A wins, B wins, or confirm no-contest), preventing abuse.

% =====================================================================
\section{The VIZ Token in the Prediction Market Experiment}\label{sec:token}

\subsection{Token Utility}

The VIZ token serves four distinct functions within the Onix Protocol ecosystem:

\begin{enumerate}
  \item \textbf{Medium of exchange.} All bets, LP deposits, oracle insurance,
  and dispute fees are denominated in VIZ. The protocol never mints tokens --- it
  is strictly zero-sum at the consensus level.
  \item \textbf{Governance weight.} VIZ staked as vesting shares (or deposited in
  the Lazy Pool) confers voting power in DAO committee disputes and chain
  parameter governance. This creates direct utility for token holding beyond
  speculation.
  \item \textbf{Oracle collateral.} Oracle insurance bonds are denominated in
  VIZ. The bond must exceed the oracle's potential manipulation profit, creating
  a demand for token accumulation by oracle operators.
  \item \textbf{Dispute escrow.} Dispute fees are escrowed in VIZ, creating a
  cost for frivolous disputes and a reward mechanism for valid challenges.
\end{enumerate}

\subsection{Token Demand Feedback Loop}

The protocol creates a structural demand cycle:

\begin{align*}
\text{LP deposits} &\xrightarrow{\text{lock}} \text{reduced circulating supply} \\
\text{Oracle insurance} &\xrightarrow{\text{lock}} \text{reduced circulating supply} \\
\text{Active bets} &\xrightarrow{\text{lock}} \text{reduced circulating supply} \\
\text{Governance weight} &\xrightarrow{\text{utility}} \text{demand for staking}
\end{align*}

The total locked supply at any time is:

\[
S_{\text{locked}} = S_{\text{pool}} + \sum_o I_o + \sum_{\text{active bets}} a_i
+ \sum_{\text{disputes}} d_k
\]

where $S_{\text{pool}}$ is the total withdrawal-locked Lazy-Pool balance (free
\emph{and} market-allocated portions are both locked for the deposit lock
period), $I_o$ is oracle insurance, $a_i$ are active bet stakes, and $d_k$ are
escrowed dispute fees. This locked supply reduces available circulating supply,
potentially creating upward price pressure as protocol usage grows --- the same
bet every protocol-native token makes, here explicitly tied to measurable
utility.

\subsection{Experimental Value of the Token}

The VIZ token in this experiment serves as a \textbf{measurement instrument}:

\begin{itemize}
  \item \textbf{Price as signal.} Token price movements in response to protocol
  events (market launches, high-volume resolutions, dispute outcomes) provide a
  continuous market-based assessment of the protocol's perceived value.
  \item \textbf{Governance participation rate.} The fraction of token holders
  participating in dispute votes measures the viability of DAO-based arbitration.
  \item \textbf{Lazy Pool accumulation rate.} The rate of pool deposits measures
  retail demand for risk-free LP yield, directly testing hypothesis Q3.
\end{itemize}

% =====================================================================
\section{Full Market Cycle: A Worked Example}\label{sec:example}

We trace a complete binary market lifecycle to illustrate the protocol's
operation.

\subsection{Setup}

\begin{itemize}
  \item \textbf{Market:} ``Will event X occur by date Y?'' (Binary: Yes/No)
  \item \textbf{Oracle:} $o$ with insurance $I_o = 5000$ VIZ
  \item \textbf{Seed liquidity:} $L = 200$ VIZ from creator
  \item \textbf{Fee parameters:} $\theta_{\text{oracle}} = 50$ bp,
  $\theta_{\text{creator}} = 50$ bp, $\theta_{\text{liq}} = 100$ bp
  \item \textbf{Lazy Pool:} $B_{\text{free}} = 10{,}000$ VIZ, $\alpha = 2\%$,
  allocates $a_{\text{alloc}} = 200$ VIZ
\end{itemize}

\begin{remark}[Units]
Amounts are displayed in VIZ for readability, but the protocol stores and
computes them as \textbf{integer milli-VIZ (mVIZ)}, with $1\text{ VIZ} =
1000\text{ mVIZ}$, and every division is floored (cf. the
Theorem~\ref{thm:lp} remark). The fee lines below are shown in mVIZ so that
$\lfloor\cdot\rfloor$ is exact; the VIZ equivalent is given alongside. This
resolves the apparent rounding mismatch that arises if one floors over
whole-VIZ quantities (e.g. $\lfloor 80\times 50/10000\rfloor$ is $0$ in VIZ but
$400$ mVIZ $=0.4$ VIZ in the real units).
\end{remark}

\subsection{Initialization}

\[
R_A = 100, \quad R_B = 100, \quad k = 10{,}000
\]

Lazy Pool auto-allocates 200 VIZ as additional LP:

\[
R_A = 200, \quad R_B = 200, \quad k = 40{,}000
\]

\subsection{Betting Phase}

\begin{itemize}
  \item \textbf{Alice} bets $a_1 = 50$ VIZ on Yes (side A):
  \[
  R'_A = 250, \quad R'_B = \lfloor 40000/250 \rfloor = 160, \quad w_1 = 200 -
  160 = 40
  \]
  \item \textbf{Bob} bets $a_2 = 80$ VIZ on No (side B):
  \[
  R'_B = 240, \quad R'_A = \lfloor 40000/240 \rfloor = 166, \quad w_2 = 250 -
  166 = 84
  \]
\end{itemize}

Implied probability: $P(\text{Yes}) = 166 / 406 \approx 41\%$, $P(\text{No}) =
240 / 406 \approx 59\%$.

\subsection{Resolution}

Oracle declares \textbf{Yes} wins. Alice is the sole winner; Bob forfeits 80
VIZ.

\begin{align*}
S_{\text{lose}} &= 80{,}000 \text{ mVIZ} \;(80 \text{ VIZ}) \\
f_{\text{oracle}} &= \lfloor 80{,}000 \times 50 / 10000 \rfloor = 400 \text{ mVIZ} \;(0.4 \text{ VIZ}) \\
f_{\text{creator}} &= \lfloor 80{,}000 \times 50 / 10000 \rfloor = 400 \text{ mVIZ} \;(0.4 \text{ VIZ}) \\
f_{\text{liq}} &= \lfloor 80{,}000 \times 100 / 10000 \rfloor = 800 \text{ mVIZ} \;(0.8 \text{ VIZ}) \\
W &= 80{,}000 - 400 - 400 - 800 = 78{,}400 \text{ mVIZ} \;(78.4 \text{ VIZ})
\end{align*}

\textbf{Alice's payout} (sole winner, $w_1 / w_1 = 1$):

\[
\pi_1 = 50{,}000 + \lfloor 78{,}400 \times 40/40 \rfloor - \tau_1 = 128{,}400
\text{ mVIZ} - \tau_1 \;(128.4 \text{ VIZ} - \tau_1)
\]

\textbf{LP return:}

\begin{itemize}
  \item Creator LP: $200 \text{ VIZ principal} + \text{time-weighted share of }
  0.8 \text{ VIZ fee pool}$
  \item Lazy Pool LP: $200 \text{ VIZ principal} + \text{time-weighted share of }
  0.8 \text{ VIZ fee pool}$
\end{itemize}

\textbf{Verification} (Theorem~\ref{thm:lp}):

\begin{align*}
\text{Money}_{\text{IN}} &= 200 + 200 + 50 + 80 = 530 \\
\text{Money}_{\text{OUT}} &= 200 + 200 + 50 + 78.4 + 0.4 + 0.4 + 0.8 = 530 \quad
\checkmark
\end{align*}

\subsection{Dispute Scenario}

If Bob disputes within $\Delta t_{\text{grace}} = 12\text{h}$, paying $d = 10$
VIZ:

\begin{enumerate}
  \item All payouts freeze.
  \item Oracle must respond within 12 hours.
  \item In committee mode, all token holders vote (stake-weighted, public,
  revisable).
  \item If dispute upheld: payouts recalculated with corrected outcome, oracle
  insurance slashed.
  \item If dispute rejected: original payouts proceed, Bob loses 10 VIZ (split
  50/50 to oracle and resolver).
  \item If no resolution within 14 days: auto-close with full refunds.
\end{enumerate}

\subsection{Early-Exit Scenario}\label{sec:exitexample}

Suppose that instead of holding, Alice cancels her Yes bet mid-market after
Bob's bet has moved the curve in her favor, so the curve refund exceeds her 50
VIZ stake. Under the deferred-claim mechanism (\Cref{sec:earlyexit}) she is
returned her principal immediately, $\text{refund} = \min(\text{curve\_refund},
50) = 50$ VIZ, and the profit tail $\max(\text{curve\_refund} - 50, 0)$ is
recorded as a deferred claim on Yes rather than paid from the curve. At
resolution, if Yes wins, her claim is paid FIFO from the capped bucket
$\lfloor 0.33 \times S_{\text{lose}} \rfloor$; if that bucket is already drained
by earlier exiters, she takes a haircut; if No wins, her claim pays zero. LP
principal is untouched in every branch, and the winners who \emph{held} to
resolution receive at least $(1 - 0.33)\,S_{\text{lose}}$ of the losing pool plus
any unused bucket --- concretely illustrating Theorem~\ref{thm:earlyexit}.

% =====================================================================
\section{Anti-MEV: Batch and Commit-Reveal Betting}\label{sec:antimev}

\subsection{MEV in Prediction Markets}

In continuous-time prediction markets, front-running and sandwich attacks
extract value from bettors. When a large bet is broadcast to the mempool, an
attacker can:

\begin{enumerate}
  \item Front-run: place a bet on the same side before the large bet, profiting
  from the price movement
  \item Sandwich: place bets on both sides around the large bet
\end{enumerate}

\subsection{Batch Settlement}

The protocol supports opt-in batch betting (binary markets). Bets submitted
within an epoch of $E$ blocks are queued and settled at a \textbf{uniform price}
at the epoch boundary. Only the net residual (aggregate demand difference) moves
the AMM:

\[
\Delta R_A = \sum_{\text{batch}} a_{i,A} - \sum_{\text{batch}} a_{i,B}
\]

This eliminates intra-batch ordering advantage: all bets in a batch receive
identical pricing regardless of submission order.

\subsection{Commit-Reveal}

For stronger MEV protection, bettors may use a two-phase commit-reveal scheme:

\begin{enumerate}
  \item \textbf{Commit:} Submit $H(\text{bet} \| \text{nonce})$ with escrow. The
  bet's direction and amount are hidden.
  \item \textbf{Reveal:} After the epoch closes, submit $(\text{bet},
  \text{nonce})$ to execute at the batch price.
\end{enumerate}

Unrevealed commitments forfeit a governance-defined penalty percentage of the
escrow, preventing spam commitments.

% =====================================================================
\section{Implementation}\label{sec:impl}

\subsection{Consensus-Level Operations on VIZ DLT}

The protocol is implemented as first-class consensus-validated operations on the
VIZ distributed ledger --- not smart contracts, not custom payloads. VIZ DLT
provides ${\sim}3$-second block times, Delegated Proof of Stake consensus, and
named accounts (Graphene-style) with no general-purpose virtual machine.

Every financial action (market creation, bet placement, oracle resolution,
dispute filing, LP deposit/withdrawal) is a \texttt{pm\_*} operation validated by
every validator node. Invalid operations are rejected before block inclusion.

\begin{center}
\begin{tabular}{lll}
\hline
\textbf{Layer} & \textbf{Examples} & \textbf{Consensus-Validated} \\
\hline
Protocol operations & \texttt{pm\_create\_market}, \texttt{pm\_place\_bet}, \ldots & Yes --- every node validates \\
Virtual operations & \texttt{pm\_payout}, \texttt{pm\_early\_exit\_claim\_paid}, \ldots & Yes --- deterministic, block-time \\
Metadata & Market descriptions, dispute evidence & No --- client-side only \\
\hline
\end{tabular}
\end{center}

\subsection{On-Chain State}\label{sec:onchainstate}

All protocol state resides in chainbase indexed objects. Key objects include:

\begin{itemize}
  \item \texttt{pm\_market\_object}: market configuration, CPMM reserves, fee parameters, state
  \item \texttt{pm\_bet\_object}: bettor positions with weight and time penalty
  \item \texttt{pm\_liquidity\_object}: LP positions with time-weight
  \item \texttt{pm\_lazy\_pool\_object}: singleton pool state (balances, shares, reward accumulator)
  \item \texttt{pm\_oracle\_object}: oracle registration, insurance, reputation counters
  \item \texttt{pm\_dispute\_object}: dispute state and resolution
  \item \texttt{pm\_deferred\_claim\_object}: early-exit outcome-contingent claims (FIFO by exit time)
\end{itemize}

Reputation metrics are computed on read (not stored), ensuring consistency
without additional write operations.

\subsection{Governance Parameters}\label{sec:governance}

All economic parameters (fees, penalties, insurance requirements, dispute
windows, lazy pool settings, and the early-exit reward cap
\texttt{pm\_early\_exit\_reward\_cap\_percent}) are delegate median-voted. Each
elected validator publishes preferred values; the network computes the median.
Parameters change without hard forks or deployments. Kill-switches allow
governance to disable subsystems (leverage, commit-reveal) by median vote.

\subsection{Scalability}\label{sec:scalability}

Because every operation is consensus-validated, the protocol's cost model
matters at scale (thousands of concurrent markets and bets). The design keeps
per-block and per-operation work bounded:

\begin{itemize}
  \item \textbf{$O(1)$ user-facing operations.} Bet placement, cancellation, and
  LP deposit/withdraw touch a fixed number of indexed objects (the market, the
  bettor's position, and --- for the pool --- the singleton accumulator). None
  iterates over all participants. Reward distribution uses the MasterChef
  accumulator $\rho$ (\Cref{sec:lazyacct}), so a pool with $n$ depositors settles
  rewards in $O(1)$ per actor rather than $O(n)$ per resolution.
  \item \textbf{Bounded deferred work per block.} Settlement is the only fan-out
  step (a market with $m$ winning bets generates $m$ \texttt{pm\_payout} virtual
  operations; early-exit distribution adds at most one
  \texttt{pm\_early\_exit\_claim\_paid} per surviving deferred claim, itself
  capped per market). To prevent a single large resolution from bloating a
  block, payout and cron processing are rate-limited by the median parameter
  \texttt{pm\_processing\_cap\_per\_block}: at most a fixed number of
  payouts/cron items are processed per block, and the remainder carry to
  subsequent blocks. A round-robin cursor over markets ensures new markets are
  not starved by low-id incumbents. Worst-case per-block work is therefore
  $O(\text{cap})$, independent of how many markets resolve in the same interval.
  \item \textbf{No write amplification from reputation.} The 14 oracle metrics
  are counters updated only on the oracle's own actions; the composite
  reliability score is computed \textbf{on read}, not written each block
  (\Cref{sec:onchainstate}). Market discovery metadata is built by a
  non-consensus plugin and never enters block validation.
  \item \textbf{State growth.} State is linear in live objects (markets, open
  positions, active disputes, unsettled deferred claims). Resolved markets and
  paid positions are terminal and prunable by clients; consensus retains only
  what open obligations require.
\end{itemize}

The binding constraint at very high market counts is thus block space for
settlement fan-out, which \texttt{pm\_processing\_cap\_per\_block} converts from
a latency spike into bounded, amortized throughput rather than a
consensus-halting cost. A quantitative stress simulation across many concurrent
markets of varying volume and volatility is part of the experimental program
(\Cref{sec:hypotheses}, H4).

% =====================================================================
\section{Discussion}\label{sec:discussion}

\subsection{Comparison with Existing Approaches}

\begin{center}
\begin{tabular}{lllll}
\hline
\textbf{Dimension} & \textbf{Onix} & \textbf{Polymarket} & \textbf{Std LMSR} & \textbf{Uniswap AMM} \\
\hline
LP risk & \textbf{Zero} (structural) & Inventory risk & Up to $b \ln N$ & Impermanent loss \\
LP knowledge & Low (deposit) & High (manage orders) & Medium & Medium-High \\
Pricing continuity & Continuous & Discrete (book) & Continuous & Continuous \\
Fee extraction & Losers only, at resolution & Bid-ask spread & Spread & Every trade \\
Oracle model & Bonded + DAO dispute & UMA Oracle & Operator & N/A \\
Price coherence & Math invariant & Arbitrage-dependent & Math invariant & Math invariant \\
CTF split/merge & No & Yes & No & No \\
\hline
\end{tabular}
\end{center}

\subsection{Limitations and Honest Tradeoffs}\label{sec:limitations}

\begin{enumerate}
  \item \textbf{LP yield is volume-dependent, not depth-dependent.} A market with
  high subsidy and low volume earns the same absolute fees as one with low
  subsidy and equal volume. The subsidy provides depth (lower slippage) but not
  yield.
  \item \textbf{LP profit is not guaranteed.} If a market resolves with zero
  losing bets, no fees are generated. LP principal is returned, but yield may be
  zero.
  \item \textbf{Parimutuel tokens are not fixed-value instruments.} Unlike
  standard LMSR where 1 winning token $=$ 1 currency unit, Onix tokens are
  proportional claims on the losers' pool. If all bettors pick the winner,
  everyone breaks even.
  \item \textbf{DPoS governance tradeoffs.} The delegate-voted parameter model
  has known concentration risks inherent to DPoS (shared with EOS, Hive, Tron).
  The Lazy Pool governance-weight mechanism partially mitigates this by
  broadening the effective electorate.
  \item \textbf{Token liquidity dependency.} Economic guarantees (insurance
  bonds, dispute fees) scale with token market value. The protocol assumes
  utility drives demand --- the standard assumption for protocol-native tokens.
  \item \textbf{Leverage adds bounded credit risk to the pool.} The opt-in
  leverage subsystem (\Cref{sec:leverage}) is the one component \emph{not}
  covered by Theorem~\ref{thm:lp}. When enabled, the pool acts as a lender and
  can incur a shortfall on the same-side-cancel path, bounded by borrower
  collateral and isolated to a governed fund fraction. We claim full loan
  recovery on the cascade and settlement paths as a \emph{design property
  verified in simulation}, not as a proven theorem; a closed-form recovery proof
  under arbitrary reserve trajectories is open work. The subsystem is default-off
  and disableable by median vote, so the unconditional LP guarantee can always
  be restored.
  \item \textbf{Early exit is a bounded discount, and the cap is a policy
  choice.} The deferred-claim mechanism (\Cref{sec:earlyexit}) gives a
  \emph{hard}, unconditional LP guarantee (Theorem~\ref{thm:earlyexit}, via the
  headroom clamp), but at the cost of making early exit a capped, FIFO-ordered,
  outcome-contingent payoff rather than a mark-to-curve cash-out. The cap $c$ is a
  governance policy choice: too low a cap makes early exit unattractive and thins
  secondary liquidity, too high a cap erodes the reward to holders --- and in the
  high-fee corner the settlement clamp further shrinks the effective bucket to
  preserve solvency, deepening the early-exit haircut. Finding the value that
  maximizes secondary liquidity without discouraging holding is an empirical
  question (\Cref{sec:hypotheses}, H8), not one the theory settles.
  \item \textbf{Public dispute voting has a timing tradeoff.} Open, revisable
  committee ballots (\Cref{sec:disputemodes}) admit a last-mover advantage and
  were chosen deliberately over commit-reveal for auditability. This is a value
  judgment (transparency over secrecy), not a proof that open voting is
  manipulation-optimal; deployments that prioritize secrecy would need
  account-mode resolution instead.
\end{enumerate}

\subsection{Experimental Hypotheses}\label{sec:hypotheses}

The deployed protocol is designed to test the following hypotheses:

\textbf{H1 (Liquidity flywheel).} \emph{Zero-risk LP provision attracts retail
capital sufficient to produce market depth that meaningfully reduces slippage
relative to comparable platforms.} Measurable: Lazy Pool total deposits, average
market depth (reserve ratios), slippage per unit bet size.

\textbf{H2 (Information aggregation).} \emph{AMM/LMSR pricing with parimutuel
settlement produces probability estimates of comparable accuracy to CLOB-based
prediction markets for equivalent information conditions.} Measurable: Brier
scores, calibration curves, comparison with external benchmarks.

\textbf{H3 (Oracle reliability under DAO governance).} \emph{Bonded oracles with
committee-mode dispute resolution achieve resolution accuracy comparable to
centralized oracle services, with dispute rates below a sustainable threshold.}
Measurable: oracle dispute rate, dispute loss rate, average resolution time,
reliability score distribution.

\textbf{H4 (Lazy Pool sustainability).} \emph{The graduated recall and fault
stamp mechanisms prevent oracle exploitation of the Lazy Pool, maintaining
positive net yield for pool depositors across diverse market portfolios.}
Measurable: pool NAV over time, recall frequency, fault stamp distribution, net
yield per share.

\textbf{H5 (Token demand correlation).} \emph{Protocol usage (volume, market
count, LP deposits) positively correlates with token staking demand and
governance participation rate.} Measurable: token staking rate, dispute vote
participation, Lazy Pool deposit rate, correlation analysis.

\textbf{H6 (Anti-MEV effectiveness).} \emph{Batch settlement and commit-reveal
mechanisms reduce measurable MEV extraction compared to continuous instant
betting.} Measurable: price impact asymmetry (batch vs. instant), sandwich
attack frequency, bettor execution quality.

\textbf{H7 (Leverage solvency).} \emph{Under live trajectories, snapshot-based
liquidation recovers the loan on the cascade and settlement paths, confining
realized pool shortfall to the same-side-cancel path and within borrower
collateral, so that leverage interest is net-accretive to pool yield.} (This
tests the design property of \Cref{sec:leverage} empirically, in lieu of a
closed-form recovery proof.) Measurable: realized loan-recovery ratio per
liquidation event ($\text{recovered}/\lambda m$), frequency and magnitude of
same-side-cancel shortfalls relative to posted collateral, leverage interest as a
share of total pool yield, pool NAV with vs. without leverage enabled.

\textbf{H8 (Early-exit liquidity discount).} \emph{The deferred-claim mechanism
(\Cref{sec:earlyexit}) provides useful secondary liquidity --- bettors and
leveraged positions do exit early --- while the capped, FIFO bucket keeps the
early-exit payoff strictly below the hold payoff, so held winners are not
diluted and no early exit ever draws on LP principal.} Measurable: early-exit
volume and frequency, distribution of \texttt{paid}/\texttt{claimed} (haircut)
ratios from the \texttt{pm\_early\_exit\_claim\_paid} virtual operation, realized
early-exit yield vs. hold yield on matched positions, and the empirical
$\text{uncovered}$ rate (expected: identically zero) as a live check of
Theorem~\ref{thm:earlyexit}.

% =====================================================================
\section{Conclusion}\label{sec:conclusion}

We have presented the Onix Protocol, a prediction market architecture that
achieves a structural LP principal guarantee by decoupling the pricing function
(CPMM for binary, LMSR for multi-outcome) from the settlement function
(parimutuel, losers-fund-winners). We proved that this guarantee holds for both
market types regardless of weight assignment, fee parameters, or betting
distribution (Theorem~\ref{thm:lp}), and that it \emph{survives the presence of a
curve that also prices early exits} --- the deferred outcome-contingent claim
mechanism (Theorem~\ref{thm:earlyexit}) lets bettors and leveraged positions
leave early without ever realizing a trading P\&L against LP depth, with no
uncovered shortfall and no minting.

The ``Lazy'' Pool mechanism enables passive retail LP participation with
automated capital deployment and graduated opportunity-cost protection. The
two-mode dispute system --- combining bonded oracles with stake-weighted DAO
committee voting --- provides a viable alternative to centralized arbitration
while maintaining oracle incentive compatibility.

The protocol is implemented as consensus-level operations on the VIZ distributed
ledger and is designed as an experiment testing eight explicit hypotheses about
liquidity bootstrapping, information aggregation accuracy, DAO governance
viability, leverage solvency, early-exit liquidity, and token demand dynamics.
The opt-in leverage subsystem is the single component that trades the
unconditional LP guarantee for bounded, default-off pool credit risk; everything
else preserves the structural guarantees of Theorems~\ref{thm:lp}
and~\ref{thm:earlyexit}.

The mathematical properties are verifiable. The economic hypotheses will be
tested by market participation.

% =====================================================================
\begin{thebibliography}{99}

\bibitem{hanson2003} Hanson, R. (2003). Combinatorial Information Market Design.
\emph{Information Systems Frontiers}, 5(1), 107--119.
\url{https://doi.org/10.1023/A:1022058209073}

\bibitem{uniswapv2} Adams, H., Zinsmeister, N., \& Robinson, D. (2020).
\emph{Uniswap v2 Core.} Uniswap Labs. \url{https://uniswap.org/whitepaper.pdf}

\bibitem{uniswapv3} Adams, H., Zinsmeister, N., Salem, M., Keefer, R., \&
Robinson, D. (2021). \emph{Uniswap v3 Core.} Uniswap Labs.
\url{https://uniswap.org/whitepaper-v3.pdf}

\bibitem{berg2008} Berg, J., Forsythe, R., Nelson, F., \& Rietz, T. (2008).
Results from a Dozen Years of Election Futures Markets Research. In
\emph{Handbook of Experimental Economics Results} (Vol. 1, pp. 742--751).
Elsevier. \url{https://doi.org/10.1016/S1574-0722(07)00080-7}

\bibitem{cowgill2009} Cowgill, B., Wolfers, J., \& Zitzewitz, E. (2009). Using
Prediction Markets to Track Information Flows: Evidence from Google. In
\emph{Auctions, Market Mechanisms and Their Applications (AMMA 2009)}, LNICST
Vol. 14. Springer. \url{https://doi.org/10.1007/978-3-642-03821-1_2}

\bibitem{levitt2004} Levitt, S. D. (2004). Why Are Gambling Markets Organised So
Differently from Financial Markets? \emph{The Economic Journal}, 114(495),
223--246. \url{https://doi.org/10.1111/j.1468-0297.2004.00207.x}

\bibitem{gnosisctf} Gnosis. \emph{Conditional Tokens Framework (CTF)
Documentation.} \url{https://docs.gnosis.io/conditionaltokens/}

\bibitem{uma} UMA Protocol. \emph{Optimistic Oracle Documentation.}
\url{https://docs.uma.xyz/}

\bibitem{compound} Leshner, R., \& Hayes, G. (2019). \emph{Compound: The Money
Market Protocol.} Compound Labs.
\url{https://compound.finance/documents/Compound.Whitepaper.pdf}

\bibitem{masterchef} SushiSwap. (2020). \emph{MasterChef Contract.}
\url{https://github.com/sushiswap/sushiswap}

\bibitem{viz} Piskunov, A. \emph{VIZ: Distributed Ledger Technical Description,
Fair DPoS, and Governance.} VIZ-Blockchain.
\url{https://viz-blockchain.github.io/viz-cpp-node/}

\bibitem{arrow2008} Arrow, K. J., Forsythe, R., Gorham, M., Hahn, R., Hanson,
R., Ledyard, J. O., et al. (2008). The Promise of Prediction Markets.
\emph{Science}, 320(5878), 877--878. \url{https://doi.org/10.1126/science.1157679}

\bibitem{wolfers2004} Wolfers, J., \& Zitzewitz, E. (2004). Prediction Markets in
Theory and Practice. \emph{NBER Working Paper} No. 10248 (also \emph{Journal of
Economic Perspectives}, 18(2), 107--126). \url{https://doi.org/10.3386/w10248}

\bibitem{othman2013} Othman, A., Pennock, D. M., Reeves, D. M., \& Sandholm, T.
(2013). A Practical Liquidity-Sensitive Automated Market Maker. \emph{ACM
Transactions on Economics and Computation}, 1(3), Article 14.
\url{https://doi.org/10.1145/2509413.2509414}

\bibitem{abernethy2011} Abernethy, J., Chen, Y., \& Vaughan, J. W. (2011). An
Optimization-Based Framework for Automated Market-Making. In \emph{Proceedings of
the 12th ACM Conference on Electronic Commerce (EC '11)} (pp. 297--306).
\url{https://doi.org/10.1145/1993574.1993621}

\end{thebibliography}

\end{document}
